% \subsection{Brain--Computer Interfaces}

The BCI Society adopted a working definition of a brain--computer interface (BCI) as \emph{``a system that measures brain activity and converts it in (nearly) real-time into functionally useful outputs to replace, restore, enhance, supplement, and/or improve the natural outputs of the brain, thereby changing the ongoing interactions between the brain and its external or internal environments. It may additionally modify brain activity using targeted delivery of stimuli to create functionally useful inputs to the brain.''}~\cite{bci-soc}

In practical terms, a BCI provides a direct communication pathway between the brain and an external device, translating the user’s intentions---as reflected in their brain signals---into control commands without relying on normal neuromuscular pathways. These systems can be used to replace or restore lost functions, such as enabling communication for individuals with severe motor impairments, or to enhance and supplement existing functions by allowing the brain to directly interact with external software or hardware~\cite{bci-soc}.

A typical BCI system operates as a processing pipeline composed of several stages. While implementation details vary across modalities and applications, these systems include signal acquisition, signal processing, command generation, and an application interface with feedback~\cite{wang-2025}.

\paragraph*{Signal Acquisition}
The signal acquisition stage involves capturing raw brain signals through a hardware--software interface. Brain activity can be recorded using non-invasive techniques (e.g., \gls{eeg}) or invasive approaches (e.g., electrocorticography), depending on application requirements and performance constraints. Among these, \gls{eeg} remains the most commonly used modality in BCI research due to its non-invasiveness and high temporal resolution. Regardless of the recording technique, the primary objective of this stage is to acquire neural data and stream it to subsequent processing steps~\cite{wang-2025}.

\paragraph*{Signal Processing}
Once acquired, brain signals are processed to extract information relevant to the user’s intention or cognitive state. This stage typically includes preprocessing to reduce noise and artifacts, followed by feature extraction and classification. From a signal processing perspective, the extracted features and classification strategies are closely linked to the interaction paradigm adopted.

At a high level, these paradigms can be divided into passive and active categories. Passive paradigms are primarily used for state monitoring and do not aim to decode voluntary user commands. Instead, signal processing focuses on identifying neural patterns associated with involuntary or spontaneous states such as attention, mental workload, fatigue, or emotional state. The information derived from these patterns is typically used to adapt system behavior rather than to generate explicit control commands.

Active paradigms, in contrast, aim to decode voluntary and intentional user actions. In this case, signal processing is designed to extract features that explicitly encode the user’s desired command. Active paradigms encompass a range of approaches, including motor imagery as well as stimulus-based paradigms relying on visual, auditory, tactile, or olfactory modalities. Within this category, a further conceptual distinction can be made between endogenous and exogenous mechanisms. Endogenous paradigms rely on self-generated neural activity and do not require external stimulation; motor imagery is a canonical example, where imagined movements induce modulations in sensorimotor rhythms, such as event-related desynchronization or synchronization in the mu and beta bands. Exogenous paradigms, by contrast, are stimulus-driven and exploit neural responses elicited by external inputs. For instance, a flickering visual element in an interface can evoke steady-state visually evoked potentials (SSVEP), characterized by oscillatory brain responses at the stimulation frequency, while infrequent or task-relevant stimuli can elicit event-related potentials such as the P300. In these cases, user intent is inferred by analyzing the neural response to the presented stimulus. 

Hybrid paradigms may combine multiple mechanisms, for example passive attention features with SSVEP-based selection, in order to improve robustness, increase the number of available control commands, or mitigate limitations of individual paradigms~\cite{wang-2025}.

\paragraph*{Command Generation}
Following feature extraction and classification, the command generation stage maps the recognized brain patterns to device-specific control signals. Depending on the application, this may involve selecting symbols, triggering discrete actions, or continuously controlling external devices such as cursors, robots, or assistive technologies.

\paragraph*{Application and Feedback}
The final stage delivers the generated command to the external system and provides feedback to the user. Feedback may be visual, auditory, sensorial, or multimodal, and plays a central role in supporting user learning and adaptation. Most BCI systems operate in a closed-loop configuration, in which feedback allows the user to adjust their mental strategy and contributes to improved system performance over time~\cite{wang-2025}.
